\textbf{Phần 1. Trắc nghiệm lựa chọn}
\setcounterpageref{ex}{0}

\begin{ex} %[3H5N33]
Đơn vị đo nhiệt độ thường dùng ở nước ta là
\choice
{Độ F}
{Độ K}
{Độ V}
{\True Độ C ($^\circ\text{C}$)}
\loigiai{
Đơn vị đo nhiệt độ thông dụng là độ C ($^\circ\text{C}$).\\
\textbf{Đáp án của câu hỏi là D}
}
\end{ex}

\begin{ex} %[3H5N33]
Dụng cụ dùng để đo nhiệt độ là
\choice
{Cân đồng hồ}
{Thước cuộn}
{\True Nhiệt kế}
{Com-pa}
\loigiai{
Nhiệt kế dùng để đo nhiệt độ.\\
\textbf{Đáp án của câu hỏi là C}
}
\end{ex}

\begin{ex} %[3H5H33]
Nhiệt độ cơ thể người khỏe mạnh bình thường vào khoảng
\choice
{\True $37^\circ\text{C}$}
{$30^\circ\text{C}$}
{$42^\circ\text{C}$}
{$100^\circ\text{C}$}
\loigiai{
Thân nhiệt bình thường của người khoảng $37^\circ\text{C}$.\\
\textbf{Đáp án của câu hỏi là A}
}
\end{ex}

\begin{ex} %[3H5H33]
Nước sôi ở nhiệt độ bao nhiêu độ C?
\choice
{$0^\circ\text{C}$}
{$50^\circ\text{C}$}
{\True $100^\circ\text{C}$}
{$37^\circ\text{C}$}
\loigiai{
Nước sôi ở $100^\circ\text{C}$.\\
\textbf{Đáp án của câu hỏi là C}
}
\end{ex}

\begin{ex} %[3H5H33]
Nước đóng băng (nước đá) ở nhiệt độ bao nhiêu độ C?
\choice
{$10^\circ\text{C}$}
{\True $0^\circ\text{C}$}
{$100^\circ\text{C}$}
{$-50^\circ\text{C}$}
\loigiai{
Nước đóng băng ở $0^\circ\text{C}$.\\
\textbf{Đáp án của câu hỏi là B}
}
\end{ex}

\begin{ex} %[3H5H33]
Nhiệt độ buổi sáng là $22^\circ\text{C}$, buổi trưa nhiệt độ tăng thêm $8^\circ\text{C}$. Nhiệt độ buổi trưa là
\choice
{\True $30^\circ\text{C}$}
{$28^\circ\text{C}$}
{$32^\circ\text{C}$}
{$35^\circ\text{C}$}
\loigiai{
$22 + 8 = 30^\circ\text{C}$.\\
\textbf{Đáp án của câu hỏi là A}
}
\end{ex}

\begin{ex} %[3H5V33]
Nhiệt độ hôm nay ở Hà Nội là $32^\circ\text{C}$, ở Sa Pa là $18^\circ\text{C}$. Nhiệt độ ở Hà Nội cao hơn ở Sa Pa là
\choice
{$12^\circ\text{C}$}
{\True $14^\circ\text{C}$}
{$16^\circ\text{C}$}
{$18^\circ\text{C}$}
\loigiai{
$32 - 18 = 14^\circ\text{C}$.\\
\textbf{Đáp án của câu hỏi là B}
}
\end{ex}

\begin{ex} %[3H5V33]
Nhiệt kế đo ngoài trời chỉ $35^\circ\text{C}$, sau cơn mưa nhiệt độ giảm đi $6^\circ\text{C}$. Nhiệt độ sau mưa là
\choice
{$28^\circ\text{C}$}
{$30^\circ\text{C}$}
{$31^\circ\text{C}$}
{\True $29^\circ\text{C}$}
\loigiai{
$35 - 6 = 29^\circ\text{C}$.\\
\textbf{Đáp án của câu hỏi là D}
}
\end{ex}

\begin{ex} %[3H5V33]
Khi bạn An bị sốt, cặp nhiệt độ thấy chỉ $38{,}5^\circ\text{C}$. Nhiệt độ của An cao hơn thân nhiệt bình thường ($37^\circ\text{C}$) là
\choice
{$1^\circ\text{C}$}
{$2^\circ\text{C}$}
{\True $1{,}5^\circ\text{C}$}
{$0{,}5^\circ\text{C}$}
\loigiai{
$38{,}5 - 37 = 1{,}5^\circ\text{C}$.\\
\textbf{Đáp án của câu hỏi là C}
}
\end{ex}

\begin{ex} %[3H5V33]
Một cốc nước nóng $60^\circ\text{C}$, để nguội bớt $25^\circ\text{C}$. Nhiệt độ cốc nước lúc này là
\choice
{$30^\circ\text{C}$}
{\True $35^\circ\text{C}$}
{$40^\circ\text{C}$}
{$45^\circ\text{C}$}
\loigiai{
$60 - 25 = 35^\circ\text{C}$.\\
\textbf{Đáp án của câu hỏi là B}
}
\end{ex}

\begin{ex} %[3H5V33]
Nhiệt độ phòng lúc chưa bật điều hòa là $33^\circ\text{C}$, sau khi bật điều hòa nhiệt độ giảm xuống còn $26^\circ\text{C}$. Nhiệt độ phòng đã giảm bao nhiêu độ C?
\choice
{$6^\circ\text{C}$}
{$8^\circ\text{C}$}
{$9^\circ\text{C}$}
{\True $7^\circ\text{C}$}
\loigiai{
$33 - 26 = 7^\circ\text{C}$.\\
\textbf{Đáp án của câu hỏi là D}
}
\end{ex}

\begin{ex} %[3H5C33]
Nhiệt độ ngày thứ Hai là $28^\circ\text{C}$, ngày thứ Ba tăng thêm $3^\circ\text{C}$, ngày thứ Tư lại giảm đi $5^\circ\text{C}$ so với ngày thứ Ba. Nhiệt độ ngày thứ Tư là
\choice
{\True $26^\circ\text{C}$}
{$25^\circ\text{C}$}
{$27^\circ\text{C}$}
{$28^\circ\text{C}$}
\loigiai{
Thứ Ba: $28 + 3 = 31^\circ\text{C}$.\\
Thứ Tư: $31 - 5 = 26^\circ\text{C}$.\\
\textbf{Đáp án của câu hỏi là A}
}
\end{ex}

\textbf{Phần 2. Trắc nghiệm đúng sai}
\setcounterpageref{ex}{0}

\begin{ex} %[3H5H33]
Xét tính đúng sai về nhiệt độ:
\choiceTF
{\True Nước đá đang tan có nhiệt độ là $0^\circ\text{C}$.}
{\True Nước đang sôi có nhiệt độ là $100^\circ\text{C}$.}
{Nhiệt độ $40^\circ\text{C}$ lạnh hơn nhiệt độ $20^\circ\text{C}$.}
{\True Nhiệt độ $40^\circ\text{C}$ nóng hơn nhiệt độ $20^\circ\text{C}$.}
\loigiai{
\begin{itemize}
    \item Nước đá tan ở $0^\circ\text{C}$: Đúng.
    \item Nước sôi ở $100^\circ\text{C}$: Đúng.
    \item Nhiệt độ càng cao thì càng nóng. Ý 3 Sai, Ý 4 Đúng.
\end{itemize}
\begin{itemize}
    \item a) \textbf{Đúng}.
    \item b) \textbf{Đúng}.
    \item c) \textbf{Sai}.
    \item d) \textbf{Đúng}.
\end{itemize}
}
\end{ex}

\begin{ex} %[3H5V33]
Theo dõi nhiệt độ 4 ngày trong tuần: Thứ Hai $26^\circ\text{C}$, Thứ Ba $28^\circ\text{C}$, Thứ Tư $30^\circ\text{C}$, Thứ Năm $25^\circ\text{C}$. Xét tính đúng sai:
\choiceTF
{\True Ngày nóng nhất là Thứ Tư ($30^\circ\text{C}$).}
{\True Ngày mát nhất là Thứ Năm ($25^\circ\text{C}$).}
{Chênh lệch nhiệt độ giữa ngày nóng nhất và mát nhất là $6^\circ\text{C}$.}
{\True Chênh lệch nhiệt độ giữa ngày nóng nhất và mát nhất là $5^\circ\text{C}$ ($30 - 25 = 5$).}
\loigiai{
\begin{itemize}
    \item Nóng nhất: Thứ Tư $30^\circ\text{C}$. Đúng.
    \item Mát nhất: Thứ Năm $25^\circ\text{C}$. Đúng.
    \item Hiệu: $30 - 25 = 5^\circ\text{C} \neq 6^\circ\text{C}$. Ý 3 Sai, Ý 4 Đúng.
\end{itemize}
\begin{itemize}
    \item a) \textbf{Đúng}.
    \item b) \textbf{Đúng}.
    \item c) \textbf{Sai}.
    \item d) \textbf{Đúng}.
\end{itemize}
}
\end{ex}

\textbf{Phần 3. Trả lời ngắn}
\setcounterpageref{ex}{0}

\begin{ex} %[3H5H33]
Nhiệt độ nước đá đang tan là bao nhiêu độ C?
\shortans{0}
\loigiai{
\textbf{Đáp số:} $0$\\
$0^\circ\text{C}$.
}
\end{ex}

\begin{ex} %[3H5H33]
Nhiệt độ nước đang sôi là bao nhiêu độ C?
\shortans{100}
\loigiai{
\textbf{Đáp số:} $100$\\
$100^\circ\text{C}$.
}
\end{ex}

\begin{ex} %[3H5V33]
Nhiệt độ sáng sớm là $19^\circ\text{C}$, đến trưa tăng thêm $7^\circ\text{C}$. Hỏi nhiệt độ buổi trưa là bao nhiêu độ C?
\shortans{26}
\loigiai{
\textbf{Đáp số:} $26$\\
$19 + 7 = 26^\circ\text{C}$.
}
\end{ex}

\begin{ex} %[3H5V33]
Nhiệt độ phòng lúc trưa là $34^\circ\text{C}$, lúc đêm giảm đi $8^\circ\text{C}$. Hỏi nhiệt độ lúc đêm là bao nhiêu độ C?
\shortans{26}
\loigiai{
\textbf{Đáp số:} $26$\\
$34 - 8 = 26^\circ\text{C}$.
}
\end{ex}

\begin{ex} %[3H5V33]
Thân nhiệt của Nam là $37^\circ\text{C}$, thân nhiệt của em bé là $38^\circ\text{C}$. Hỏi thân nhiệt em bé cao hơn Nam bao nhiêu độ C?
\shortans{1}
\loigiai{
\textbf{Đáp số:} $1$\\
$38 - 37 = 1^\circ\text{C}$.
}
\end{ex}

\begin{ex} %[3H5C33]
Một cốc nước $50^\circ\text{C}$. Sau khi cho thêm đá, nhiệt độ giảm $22^\circ\text{C}$. Cốc nước đá lúc này có nhiệt độ bao nhiêu độ C?
\shortans{28}
\loigiai{
\textbf{Đáp số:} $28$\\
$50 - 22 = 28^\circ\text{C}$.
}
\end{ex}
